% Системийн бүтэц — асуултаас хариулт хүртэл (TikZ).
% Эх: локал системийн /data#tree хуудасны зураг (docs/diagrams/Системийн_бүтэц.pdf, 11/06),
% 11/09-нөөс асаасан ресепшний хүснэгт ба «хэнд, хаана хандах» давхаргыг нэмсэн.
% Байрлалыг (x,y) координатаар өгсөн: нэг хайрцгийг зөөхөд тоог нь л өөрчилнө.
\begin{tikzpicture}[
    font=\footnotesize,
    box/.style={draw, rounded corners=2pt, align=center, text width=2.3cm,
                minimum height=0.95cm, inner sep=2pt, fill=white},
    ext/.style={box, fill=black!8},
    data/.style={box, fill=black!3},
    head/.style={font=\footnotesize\bfseries, align=center},
    arr/.style={-{Stealth[length=2mm]}, semithick},
    note/.style={font=\scriptsize, inner sep=1.5pt}]

% --- баганын гарчиг
\node[head] at (0,1.0)     {1. Асуух};
\node[head] at (3.1,1.0)   {2. Чиглүүлэх};
\node[head] at (6.2,1.0)   {3. Хайх};
\node[head] at (9.3,1.0)   {4. Сонгох};
\node[head] at (13.0,1.0)  {5. Шалгаж хариулах};

% --- 1. Асуух
\node[box]  (voice) at (0,0)      {Дуугаар асуух};
\node[ext]  (stt)   at (0,-1.7)   {Яриа таних\\{\scriptsize Gemini 3.5 Transcribe}};
\node[box]  (text)  at (0,-3.4)   {Асуултын текст};
\node[box]  (typed) at (0,-5.1)   {Бичгээр асуух};

% --- 2. Чиглүүлэх
\node[box]  (clar)  at (3.1,-1.7) {Тодруулж асуух};
\node[box]  (route) at (3.1,-3.4) {Асуулт чиглүүлэгч};

% --- 3. Хайх
\node[box]  (bm25)  at (6.2,-1.7)  {BM25\\{\scriptsize үгээр хайх}};
\node[box]  (e5)    at (6.2,-3.4)  {E5\\{\scriptsize утгаар хайх}};
\node[box]  (table) at (6.2,-5.1)  {Ресепшний хүснэгт\\{\scriptsize өрөө, анги, хуваарь}};

% --- 4. Сонгох
\node[box]  (rrf)   at (9.3,-1.7) {RRF\\{\scriptsize хоёр дүнг нэгтгэх}};
\node[box]  (sel)   at (9.3,-3.4) {Хэсэг сонгох дүрэм};

% --- 5. Шалгаж хариулах
\node[box]  (refer) at (13.0,-1.7) {Хэнд, хаана\\хандахыг заах};
\node[box]  (guard) at (13.0,-3.4) {Эхийн шалгалт};
\node[box]  (ans)   at (13.0,-5.1) {Эхтэй хариулт\\{\scriptsize эхийн үгээр}};
\node[ext]  (tts)   at (13.0,-6.8) {Дуугаар унших\\{\scriptsize Edge TTS}};

% --- сумнууд
\draw[arr] (voice) -- (stt);
\draw[arr] (stt) -- (text);
\draw[arr] (typed) -- (text);
\draw[arr] (text) -- (route);
\draw[arr] (route) -- node[note, left] {тодорхойгүй} (clar);
\draw[arr] (route.east) -- ++(0.25,0) |- (bm25.west);
\draw[arr] (route) -- (e5);
\draw[arr] (route.east) -- ++(0.25,0) |- (table.west);
\draw[arr] (bm25) -- (rrf);
\draw[arr] (e5.east) -- ++(0.25,0) |- ([yshift=-2.5mm]rrf.west);
\draw[arr] (rrf) -- (sel);
\draw[arr] (rrf) -- node[note, above] {эх алга} (refer);
\draw[arr] (sel) -- (guard);
\draw[arr] (table) -- (ans);
\draw[arr] (guard) -- node[note, right] {тэнцээгүй} (refer);
\draw[arr] (guard) -- node[note, right] {тэнцсэн} (ans);
\draw[arr] (ans) -- (tts);

% --- Мэдээлэл бэлтгэх (асуултаас өмнө тусдаа хийгдэнэ)
\draw[dashed] (-1.35,-7.8) -- (14.35,-7.8);
\node[head, anchor=west] at (-1.35,-8.2) {Мэдээлэл бэлтгэх};
\node[data] (web)    at (0,-9.2)     {Сургуулийн\\вэб сайт};
\node[data] (office) at (0,-10.9)    {Сургуулийн албадын жагсаалт, хуваарь};
\node[data] (store)  at (4.2,-10.05) {Эхийн сан\\{\scriptsize огноотой хувилбар}};
\node[data] (chunk)  at (8.4,-9.2)   {Хайлтын хэсгүүд};
\node[data] (rtab)   at (8.4,-10.9)  {Ресепшний хүснэгт};
\node[data] (tree)   at (13.0,-10.05) {Мэдлэгийн мод\\{\scriptsize 11 хүрээ, 44 салаа}};

\draw[arr] (web.east) -- ++(0.6,0) |- ([yshift=2mm]store.west);
\draw[arr] (office.east) -- ++(0.6,0) |- ([yshift=-2mm]store.west);
\draw[arr] (store.east) -- ++(0.6,0) |- (chunk.west);
\draw[arr] (store.east) -- ++(0.6,0) |- (rtab.west);
\draw[arr] (chunk.east) -- ++(0.9,0) |- ([yshift=2mm]tree.west);
\draw[arr] (rtab.east) -- ++(0.9,0) |- ([yshift=-2mm]tree.west);
\end{tikzpicture}
